\documentclass[12pt,a4paper]{article}
\usepackage[utf8]{inputenc}
\usepackage{xeCJK}
\usepackage{geometry}
\usepackage{booktabs}
\usepackage{array}
\usepackage{longtable}
\usepackage{multirow}
\usepackage{graphicx}
\usepackage{fancyhdr}
\usepackage{color}
\usepackage{xcolor}
\usepackage{amsmath}
\usepackage{amsfonts}
\usepackage{amssymb}

\usepackage{pdfpages}

\geometry{left=2cm,right=2cm,top=2.5cm,bottom=2.5cm}

\setCJKmainfont{Noto Serif CJK TC}
\setCJKsansfont{Noto Sans CJK TC}
\setCJKmonofont{Noto Sans CJK TC}

\pagestyle{fancy}
\fancyhf{}
\rhead{\thepage}
\lhead{心臟內科見習學習歷程成果報告}

\title{\Large \textbf{心臟內科見習學習歷程成果報告}\\
\large Clinical Shadowing Program Assessment Report}
\author{指導醫師：楊欽文 醫師\\
見習學生：吳羿玟 醫師}
\date{見習期間：2025年10月27日 - 2025年11月14日}

\begin{document}

\maketitle

\section{執行摘要}

吳羿玟醫師於2025年10月27日至11月14日期間完成為期三週的心臟內科臨床見習計畫。本報告基於學生的每日學習記錄、既定教學計畫以及心臟內科可信賴專業活動(EPAs)評估標準，對學習成果進行綜合評估。

\section{計畫執行與實際學習比較}

\subsection{計畫達成度分析}

\begin{table}[h]
\centering
\caption{教學計畫與實際執行對照表}
\begin{tabular}{|p{3cm}|p{5cm}|p{6cm}|}
\hline
\textbf{計畫項目} & \textbf{原定安排} & \textbf{實際執行情況} \\
\hline
門診見習 & 配合指導醫師，於新竹與生醫進行 & 實際參與多次門診，觀摩多種疾病類型 \\
\hline
病房迴診 & 每日病房迴診與病患評估 & 持續參與病房迴診，接診新病患，完成病歷撰寫 \\
\hline
心導管室實習 & 觀摩基礎至進階手術 & 觀摩ASD repair、LAAO、PCI、PTA、Carotid angiography等多種術式 \\
\hline
超音波訓練 & 標準切面學習與操作 & 學習心臟超音波標準切面、血管超音波、TEE觀摩、Contrast echo \\
\hline
心電圖判讀 & 基礎心電圖判讀訓練 & 系統性學習AV block、ST變化、心律不整等心電圖判讀 \\
\hline
特殊檢查 & 運動心電圖觀摩 & 完成運動心電圖檢查觀摩與學習 \\
\hline
\end{tabular}
\end{table}

\subsection{量化學習成果}

\begin{table}[h]
\centering
\caption{見習期間量化指標統計}
\begin{tabular}{|p{4cm}|p{2cm}|p{6cm}|}
\hline
\textbf{學習項目} & \textbf{完成數量} & \textbf{說明} \\
\hline
門診病患診療觀摩 & 50+ 例 & 涵蓋高血壓、心衰竭、冠心病、瓣膜性心臟病、PAOD等疾病 \\
\hline
病患接診 & 3 位 & 完整執行病患照護流程，包含Pericarditis等案例 \\
\hline
病歷撰寫 & 3 份 & 包括admission notes，接受指導醫師審閱修訂 \\
\hline
心臟超音波執行 & 8 次 & Apical 4-chamber view等標準切面實作 \\
\hline
血管超音波 & 5 次 & 動靜脈血管檢查技術 \\
\hline
心導管手術觀摩 & 15+ 台 & ASD repair、LAAO、PCI、PTA、FFR等多種術式 \\
\hline
TEE檢查觀摩 & 2 次 & 經食道超音波檢查，包含LAAO術前評估及IE評估 \\
\hline
Contrast echo觀摩 & 1 次 & 對比超音波檢查技術學習 \\
\hline
心電圖判讀練習 & 50+ 張 & 系統性學習各種心電圖異常 \\
\hline
運動心電圖觀摩 & 2 次 & 學習適應症、禁忌症及判讀 \\
\hline
照會案例討論 & 2 例 & Infective endocarditis等跨科合作案例 \\
\hline
\end{tabular}
\end{table}

\subsection{學習亮點}

\begin{enumerate}
\item \textbf{臨床經驗多元}：參與門診50餘例病患診療，涵蓋心臟內科常見疾病，並觀摩先天性心臟病等特殊案例
\item \textbf{實務操作紮實}：獨立接診3位病患，完成病歷撰寫，累積實際臨床經驗
\item \textbf{技能操作積極}：執行8次心臟超音波及5次血管超音波檢查，主動練習標準切面獲取
\item \textbf{手術觀摩豐富}：觀摩包括ASD repair、LAAO、複雜PCI（FFR、DK crush）、PTA等15台以上手術
\item \textbf{學習態度認真}：每日完整記錄學習心得，主動查詢資料，積極提問
\item \textbf{反思能力佳}：能夠整合臨床觀察與理論知識，提出深入思考
\item \textbf{專業態度優良}：與醫療團隊建立良好互動，尊重病患隱私
\end{enumerate}

\section{臨床技能評估}

\subsection{EPA評估結果}

根據指導醫師評估，各項EPAs達成情況如下：

\begin{longtable}{|p{1cm}|p{6cm}|p{2cm}|p{5cm}|}
\caption{心臟內科EPAs評估結果} \\
\hline
\textbf{EPA} & \textbf{專業活動} & \textbf{等級} & \textbf{評估說明} \\
\hline
\endfirsthead
\hline
\textbf{EPA} & \textbf{專業活動} & \textbf{等級} & \textbf{評估說明} \\
\hline
\endhead

EPA 1 & 心臟科病史詢問與風險評估 & Level 2 & 能系統性完成病史詢問，理解胸痛等症狀鑑別診斷 \\
\hline
EPA 2 & 心臟血管系統理學檢查 & Level 2 & 能執行基本心血管檢查，辨識心雜音，需持續練習聽診技巧 \\
\hline  
EPA 3 & 心電圖判讀與分析 & Level 2 & 能識別AV block、ST變化等常見異常，判讀能力持續進步 \\
\hline
EPA 4 & 心臟超音波基礎操作與判讀 & Level 1 & 能獲取Apical 4-chamber view等基本切面，需更多操作練習 \\
\hline
EPA 5 & 急性冠心症診斷與初步處理 & Level 1 & 理解STEMI與pericarditis鑑別，了解基本處理原則 \\
\hline
EPA 6 & 心律不整診斷與處理 & Level 1 & 能識別AFib、AV block等心律不整，了解基本治療原則 \\
\hline
EPA 7 & 心衰竭診斷與治療 & Level 1 & 能識別心衰竭症狀徵象，理解診斷標準與基本治療 \\
\hline
EPA 8 & 高血壓診斷與管理 & Level 2 & 理解血壓管理原則，知道續發性高血壓評估 \\
\hline
EPA 9 & 心導管檢查結果判讀 & Level 2 & 能判讀基本心導管影像，理解FFR、PCI等介入治療原理 \\
\hline
EPA 10 & 心臟科病患照護協調與溝通 & Level 2 & 觀察學習醫病溝通技巧，理解SDM重要性 \\
\hline
\end{longtable}

\subsection{學習成就統計}

\begin{table}[h]
\centering
\caption{EPA等級分佈統計}
\begin{tabular}{|c|c|c|}
\hline
\textbf{EPA等級} & \textbf{項目數} & \textbf{百分比} \\
\hline
Level 3 & 0 & 0\% \\
Level 2 & 6 & 60\% \\
Level 1 & 4 & 40\% \\
None & 0 & 0\% \\
\hline
\textbf{總計} & \textbf{10} & \textbf{100\%} \\
\hline
\end{tabular}
\end{table}

\section{學習品質分析}

\subsection{優勢表現}

\begin{itemize}
\item \textbf{心電圖判讀能力}：EPA 3達到Level 2，展現良好的心電圖判讀基礎，能辨識多種異常
\item \textbf{臨床觀察力}：能整合問診、理學檢查、檢查結果進行初步診斷思考
\item \textbf{學習主動性}：主動閱讀相關文獻，積極查詢不懂的醫學知識
\item \textbf{專業態度}：與醫療團隊建立良好溝通關係，展現專業素養
\item \textbf{反思能力}：每日完整記錄學習心得，能提出深入思考問題
\end{itemize}

\subsection{改進建議}

\begin{itemize}
\item \textbf{聽診技能}：需要更多實際練習以提升心雜音辨識能力，特別是女性病患聽診技巧
\item \textbf{超音波操作}：需要更多hands-on練習以提升影像獲取熟練度
\item \textbf{病歷撰寫}：可加快病歷撰寫速度，善用AI工具輔助但需保持個人風格
\item \textbf{急症處理}：建議增加急性心血管疾病模擬訓練經驗
\item \textbf{臨床決策}：建議參與更多複雜病例的醫療決策討論
\end{itemize}

\section{特色學習經驗}

\subsection{先天性心臟病}

\begin{itemize}
\item 觀摩ASD repair手術，學習肺分流量計算、repair適應症與禁忌症
\item 理解右心導管hemodynamic study的重要性
\end{itemize}

\subsection{結構性心臟病介入治療}

\begin{itemize}
\item 觀摩LAAO手術，學習CHA2DS2-VASc Score、HAS-BLED評估
\item 理解DOAC intolerance患者的治療選擇
\item 觀察術中併發症（血栓形成）的即時處理
\end{itemize}

\subsection{冠狀動脈介入治療}

\begin{itemize}
\item 學習FFR評估技術的臨床應用
\item 觀摩DK crush複雜支架技術
\item 理解塗藥氣球(DCB)的適應症與療效
\end{itemize}

\subsection{週邊血管介入治療}

\begin{itemize}
\item 觀摩洗腎廔管(AVF) PTA治療
\item 學習PAOD的Rutherford分級與治療策略
\item 理解血栓處理的技巧與風險
\end{itemize}

\subsection{新穎影像技術}

\begin{itemize}
\item 觀摩Contrast echocardiography，了解其臨床應用潛力
\item 學習TEE在瓣膜性心臟病診斷的重要性
\end{itemize}

\section{標準化評分系統}

\subsection{評分架構}

本評分系統採用多元評量方式，包含三個主要評分項目：
\begin{itemize}
\item \textbf{EPA基礎能力表現 (75\%)}：基於10項EPAs的達成等級
\item \textbf{學習態度與表現 (20\%)}：包含學習積極性、專業態度、自主學習能力
\item \textbf{特殊貢獻加分 (5\%)}：創新表現、額外學習成果、卓越表現
\end{itemize}

\subsection{EPA等級評分標準}

\begin{table}[h]
\centering
\caption{EPA等級對應分數表}
\begin{tabular}{|c|c|p{8cm}|}
\hline
\textbf{EPA等級} & \textbf{對應分數} & \textbf{能力描述} \\
\hline
Level 5 & 100分 & 完全獨立執行並能指導他人 \\
\hline
Level 4 & 98分 & 可獨立執行，監督者隨時可得 \\
\hline
Level 3 & 95分 & 可獨立執行，但需回報並接受指導 \\
\hline
Level 2 & 90分 & 可在主治醫師直接監督下執行 \\
\hline
Level 1 & 85分 & 僅能觀察學習，無法獨立操作 \\
\hline
None & 75分 & 未有學習機會（見習生合理情況）\\
\hline
\end{tabular}
\end{table}

\subsection{學習態度評分準則}

\begin{table}[h]
\centering
\caption{學習態度評分標準}
\begin{tabular}{|c|p{10cm}|}
\hline
\textbf{分數區間} & \textbf{表現標準} \\
\hline
95-100分 & 學習態度極為積極，主動尋求學習機會，完整記錄學習歷程，展現優異專業素養 \\
\hline
90-94分 & 學習態度良好，能配合教學安排，有基本學習記錄，專業表現合宜 \\
\hline
85-89分 & 學習態度一般，被動參與，學習記錄不完整 \\
\hline
80-84分 & 學習態度消極，缺乏主動性，專業表現待加強 \\
\hline
\end{tabular}
\end{table}

\subsection{吳羿玟醫師評分計算}

\subsubsection{EPA基礎能力分數}
\begin{align}
\text{EPA基礎分} &= \frac{(90 \times 6) + (85 \times 4)}{10}\\
&= \frac{540 + 340}{10}\\
&= \frac{880}{10} = 88\text{分}
\end{align}

\subsubsection{學習態度與表現評分}
根據學習記錄分析：
\begin{itemize}
\item 每日完整學習記錄（15天）：優
\item 主動查詢文獻資料：優  
\item 專業態度表現：優
\item 自我反思能力：優
\item 積極提問與討論：優
\end{itemize}
\textbf{學習態度得分：97分}

\subsubsection{特殊貢獻加分}
\begin{itemize}
\item 完整記錄三週學習歷程，展現優異自主學習能力
\item 主動研讀ISCHEMIA trial等臨床試驗
\item 學習多種先進技術（LAAO、FFR、Contrast echo、DK crush等）
\item 積極參與團隊協作與照會討論
\item 展現良好醫病溝通觀察力
\end{itemize}
\textbf{特殊貢獻得分：100分}

\subsubsection{最終總分計算}
\begin{align}
\text{總分} &= 88 \times 0.75 + 97 \times 0.20 + 100 \times 0.05\\
&= 66.0 + 19.4 + 5.0\\
&= 90.4\text{分}
\end{align}

考量見習期間的優秀學習表現、積極態度及完整學習記錄，給予指導教師裁量加分。

\textcolor{red}{\Large \textbf{最終評分：91分 (優良)}}

\subsection{評分等級對照}
\begin{table}[h]
\centering
\caption{總評成績等級表}
\begin{tabular}{|c|c|c|}
\hline
\textbf{分數區間} & \textbf{等級} & \textbf{評語} \\
\hline
90-100分 & A (優良) & 表現優異，具備良好臨床基礎 \\
\hline
80-89分 & B (良好) & 表現良好，達到預期學習目標 \\
\hline
70-79分 & C (合格) & 基本達標，仍有改進空間 \\
\hline
60-69分 & D (待改進) & 未達預期，需要額外指導 \\
\hline
60分以下 & F (不合格) & 嚴重不足，建議重新安排學習 \\
\hline
\end{tabular}
\end{table}

\section{印象深刻案例學習}

\subsection{Acute Pericarditis}

\textbf{病例}：53歲男性，急性胸痛
\begin{itemize}
\item 學習重點：Pericarditis與STEMI鑑別診斷
\item 關鍵特徵：姿勢性胸痛、Concave ST elevation、Pericardial friction rub
\item 臨床思考：ST elevation不一定是STEMI，需整合臨床表現
\end{itemize}

\subsection{Infective Endocarditis}

\textbf{病例}：53歲男性，右側感染性心內膜炎
\begin{itemize}
\item 學習重點：2023 Duke ISCVID診斷標準
\item 臨床挑戰：Substance use病史詢問技巧、病患配合度
\item 影像評估：TEE安排時機與限制
\end{itemize}

\subsection{複雜冠狀動脈介入治療}

\textbf{病例}：56歲男性，RCA狹窄接受FFR評估
\begin{itemize}
\item 學習重點：FFR在決策中的角色
\item 治療選擇：塗藥氣球 vs. 支架置放
\item 醫病溝通：如何協助病患理解複雜治療選項
\end{itemize}

\section{總結與建議}

吳羿玟醫師在三週的心臟內科見習期間展現積極的學習態度和良好的專業素養。在基礎臨床技能方面表現優良，特別是心電圖判讀和病歷撰寫能力達到良好水準。建議未來持續加強聽診技能和超音波操作能力的培養。

\subsection{未來發展建議}
\begin{enumerate}
\item 加強聽診技巧練習，特別是各種心雜音的辨識
\item 增加心臟超音波操作練習時間，建立獨立檢查能力
\item 參與臨床技能中心模擬訓練，強化操作技能
\item 持續參與心導管室實習，提升介入性治療認知
\item 參與疑難病例討論，提升臨床決策能力
\item 加強急性心血管疾病處理經驗
\end{enumerate}

\subsection{評分醫師評語- 謝慕揚}

本見習學習報告，依據吳醫師在每日見習後的記錄 (內容依照教學計畫書規範完成應記錄項目)，在見習結束後，由本人自Google Drive下載，並依據計畫書內容計分。

吳羿玟醫師在三週見習期間展現優異的學習態度，每日完整記錄學習心得，積極提問討論，主動查詢文獻資料。在臨床技能方面，心電圖判讀能力進步明顯，能辨識多種心律不整及ST-T波變化。聽診技巧仍需持續練習，特別是心雜音辨識。病歷撰寫能力良好，能在指導下完成完整病歷。

觀摩多種先進心導管技術（ASD repair、LAAO、FFR、DK crush等），展現對新技術的學習興趣。能夠思考臨床問題，提出有深度的反思。與醫療團隊互動良好，展現專業態度。

建議未來加強hands-on技能練習，特別是超音波操作。持續培養臨床思維能力，參與更多病例討論。整體表現優良，具備成為優秀臨床醫師的潛力。

\vspace{1cm}
\noindent
% \begin{tabular}{ll}
% 指導醫師簽名：& \underline{\hspace{4cm}} \\
% & 謝慕揚 醫師 \\[0.5cm]
% 指導醫師簽名：& \underline{\hspace{4cm}} \\
% & 楊欽文 醫師 \\[0.5cm]
% 日期：& \underline{\hspace{4cm}} \\
% & 2025年11月16日
% \end{tabular}

\section{學生的每日學習日誌}
\includepdf[pages=-]{Clerk 吳羿玟 2025-10-27 3 weeks daily note.pdf}

\end{document}